\documentclass[11pt,a4paper]{article}
\usepackage[T1]{fontenc}
\usepackage[utf8]{inputenc}
\usepackage{lmodern,microtype}
\usepackage[margin=25mm,top=25mm,bottom=25mm,headheight=14pt]{geometry}
\usepackage{amsmath,amssymb,amsthm,mathtools,booktabs,array,enumitem,graphicx}
\usepackage[svgnames]{xcolor}
\usepackage{fancyhdr}
\usepackage{hyperref}
\hypersetup{colorlinks=true,linkcolor=DarkSlateBlue,citecolor=DarkSlateBlue,urlcolor=DarkSlateBlue,
pdftitle={Single-point criteria for transcendence and algebraic independence},pdfsubject={A research synthesis on lacunary zeros, Mahler functions, and weighted Fibonacci sums}}
\newtheorem{theorem}{Theorem}[section]
\newtheorem{lemma}[theorem]{Lemma}
\newtheorem{proposition}[theorem]{Proposition}
\newtheorem{corollary}[theorem]{Corollary}
\theoremstyle{definition}
\newtheorem{definition}[theorem]{Definition}
\newtheorem{example}[theorem]{Example}
\theoremstyle{remark}
\newtheorem{remark}[theorem]{Remark}
\newcommand{\Q}{\mathbb Q}
\newcommand{\C}{\mathbb C}
\newcommand{\Z}{\mathbb Z}
\newcommand{\N}{\mathbb N}
\newcommand{\K}{\overline{\mathbb Q}}
\newcommand{\Qbar}{\overline{\mathbb Q}}
\DeclareMathOperator{\trdeg}{trdeg}
\DeclareMathOperator{\ord}{ord}
\DeclareMathOperator{\Res}{Res}
\newcommand{\one}{\mathbf 1}
\setlength{\emergencystretch}{2em}
\setlist{topsep=4pt,itemsep=2pt}
\pagestyle{fancy}\fancyhf{}\lhead{\small\itshape Single-point criteria for transcendence}\rhead{\small\thepage}
\renewcommand{\headrulewidth}{0.3pt}
\setcounter{tocdepth}{1}
\numberwithin{equation}{section}
\raggedbottom
\begin{document}
\begin{center}
{\Large\bfseries Single-point criteria for transcendence\\[3pt] and algebraic independence}\par
\vspace{8pt}
{\large Lacunary zeros, Mahler functions, and weighted Fibonacci sums}\par
\vspace{8pt}
{\small Research manuscript\quad $\cdot$\quad 15 November 2026}
\end{center}
\begin{abstract}
We study three ways of turning information at one point into an arithmetic conclusion. For a sufficiently sparse integer power series, irreducible truncations and a degree argument give algebraic independence of a zero and the values of all positive-order derivatives there. For Mahler functions, functional equations identify the values that must be counted together, and specialization transfers independence to algebraic arguments. For the bilateral Fibonacci sums with Chang's real cubic weight, we determine the transcendence degree of every finite subfamily. A subfamily has rank at most two if it contains no multiple-of-three index and at most three otherwise. The first two sums satisfy $T_2=T_1^2$, so together they contribute one independent generator. We give the explicit transcendence proof for Chang's common value, the joint derivative results for Stern's series and its twist, and the transcendence of specified real crossings of Lambert series and continued fractions.
\end{abstract}
\tableofcontents
\section{Arithmetic at a single point}\label{sec:transfer}
Throughout, $\K$ denotes the algebraic closure of $\Q$ in $\C$. A family is algebraically independent if every finite subfamily annihilates only the zero polynomial over the indicated ground field. Transcendence degree counts a largest algebraically independent subfamily of generators.

\begin{proposition}[Transfer through a rational expression]\label{prop:transfer}
If $t$ is transcendental over $\K$ and $R\in\K(X)$ is nonconstant, then $R(t)$ is transcendental. More generally, a nonconstant rational expression in an algebraically independent tuple is transcendental at that tuple.
\end{proposition}
\begin{proof}
Write $R=P/Q$. If the value were $c\in\K$, the tuple would annihilate $P-cQ$. This polynomial is nonzero because $R$ is nonconstant. The denominator also evaluates to a nonzero number by independence.
\end{proof}
An inverse version will be used for real crossings. Suppose $f$ is analytic on an interval $I$ and $f(\alpha)$ is transcendental at each algebraic $\alpha\in I$. Every solution in $I$ of $f(x)=c$, for a fixed $c\in\K$, is then transcendental. Existence and uniqueness of a specified crossing are separate analytic claims.

The original motivating target was Euler's constant $\gamma$. The reciprocal gamma function $G(z)=1/\Gamma(z)$ has expansion $G(z)=z+\gamma z^2+O(z^3)$ at zero, by the Euler product for $1/\Gamma(1+z)$ \cite{DLMF}. Thus $G''(0)=2\gamma$. An identity recovering a derivative supplies no arithmetic independence by itself. An application of Proposition~\ref{prop:transfer} to $\gamma$ would require an independently established transcendental parameter and an exact nonconstant rational expression for $\gamma$ in that parameter. One natural approximation uses the harmonic numbers $H_n=\sum_{k=1}^n1/k$. Integral comparison, or Euler--Maclaurin summation, gives $H_n-\log n=\gamma+1/(2n)+O(n^{-2})$. Keeping $\log n$ in the approximant supplies no integral polynomial truncation of the kind used in Section~\ref{sec:lacunary}. This approximation alone therefore does not give a lower bound by a nonzero integer resultant. The irrationality of Euler's constant remains open \cite{EulerStatus}. The following sections establish arithmetic independence for the particular functions they define.

\section{Zeros and derivatives of a lacunary series}\label{sec:lacunary}
Let $2\le N_0<N_1<\cdots$ be integers satisfying
\begin{equation}\label{eq:gaps}
 \frac{N_{n+1}}{N_n\log N_n}\longrightarrow\infty.
\end{equation}
Put $D=z\,d/dz$ and define
\[
 F(z)=1-2\sum_{j\ge0}z^{N_j},\qquad
 F_n(z)=1-2\sum_{j=0}^{n}z^{N_j}\quad (|z|<1).
\]
\begin{theorem}\label{thm:lac}
For every zero $\rho$ of $F$ in the open unit disk, the family
$\rho,F'(\rho),F''(\rho),\ldots$ is algebraically independent over $\K$.
There is exactly one positive real zero in this disk.
\end{theorem}
\begin{proof}
The series converges locally uniformly with all derivatives. On $(0,1)$ it is strictly decreasing from $1$ to $-\infty$, proving the last assertion. We prove the arithmetic assertion with $D^jF$ first. A relation with algebraic coefficients implies a nonzero relation with integer coefficients, by taking the product of its conjugates over a finite normal number field and clearing denominators.

Fix $r\ge0$ and a nonzero $P\in\Z[Z,Y_1,\ldots,Y_r]$ of total degree $b$. Suppose that $P(\rho,DF(\rho),\ldots,D^rF(\rho))=0$. Define
\[
 Q_n(z)=P(z,DF_n(z),\ldots,D^rF_n(z)).
\]
We first show $F_n$ and $Q_n$ are coprime for all sufficiently large $n$. The reciprocal polynomial $z^{N_n}F_n(1/z)$ is monic and Eisenstein at $2$, its constant coefficient being $-2$. Hence $F_n$ is irreducible over $\Q$. For $1\le j\le r$ set
\[
 W_{j,n}=\prod_{i=0}^{j-1}(D-N_{n-i})F_n .
\]
Modulo $F_n$, the polynomials $W_{1,n},\ldots,W_{r,n}$ are related to $DF_n,\ldots,D^rF_n$ by a unit triangular linear transformation. Applying its inverse in $P$ gives a nonzero polynomial $P_n^*$ of total degree at most $b$ in $Z,Y_1,\ldots,Y_r$, whose value at $z,W_{1,n},\ldots,W_{r,n}$ agrees with $Q_n$ modulo $F_n$.

Each operator kills precisely the indicated top monomial, so $\deg W_{j,n}=N_{n-j}$. For fixed $b,r$, the monomials of $P_n^*$, evaluated on $z,W_{1,n},\ldots,W_{r,n}$, have distinct leading degrees. Indeed their degrees are sums with digits between $0$ and $b$ and place values $1,N_{n-r},\ldots,N_{n-1}$. Eventually each place value exceeds $b$ times the sum of the preceding ones, by \eqref{eq:gaps}. The substitution is nonzero, and its degree is at most $b\max(1,N_{n-1})<N_n$ for large $n$. In the case $r=0$ the substitution simply gives $P(z)$. This proves the coprimality claim.

Choose $|\rho|<R<R'<1$ and a small closed disk about $\rho$ in $|z|<R$ on which $F$ has no other zero. Write $F=(z-\rho)^h U$ there with $|U|$ bounded below. For $0\le j\le r$, the tails of $D^jF_n$ on $|z|\le R$ are $O((R')^{N_{n+1}})$, since $m^j R^m/(R')^m$ decays exponentially. Rouch\'e's theorem gives a zero $\alpha_n$ of $F_n$ in the small disk for large $n$. Evaluating $F-F_n$ there shows $|\alpha_n-\rho|\ll(R')^{N_{n+1}/h}$. Local uniform boundedness of the next derivative and the assumed relation at $\rho$ now give $|Q_n(\alpha_n)|\ll e^{-c N_{n+1}}$ for a constant $c>0$.

For a polynomial $V$ let $\|V\|_1$ be its coefficient sum in absolute value. We have
\[
  \deg Q_n=O_P(N_n),\quad
  \|Q_n\|_1\le (1+(n+1)N_n^r)^{O_P(1)},\quad
  M(F_n)\le\|F_n\|_1=2n+3,
\]
where $M$ is Mahler measure; the inequality follows by averaging $\log|F_n|$ on the unit circle. For any $w$, $|Q_n(w)|\le\|Q_n\|_1\max(1,|w|)^{\deg Q_n}$. Apply this estimate at the other $N_n-1$ roots of $F_n$ in the product formula for the resultant. The leading coefficient is absorbed by $M(F_n)^{\deg Q_n}$. Since $n\le N_n$, it follows that
\[
 1\le |\Res(F_n,Q_n)|\le
 \exp\{-cN_{n+1}+O_P(N_n\log N_n)\}
 \longrightarrow0 .
\]
This contradicts the nonzero integral resultant. Finally $\rho\ne0$ and $D^j$ has highest ordinary derivative $z^j(d/dz)^j$. The invertible triangular change over $\Q(\rho)$ proves the ordinary-derivative assertion.
\end{proof}

\paragraph{The positive zero and its moments.}
For $N_0=2$, $N_{n+1}=N_n^2$ the zero is
$\rho=0.60484719466598437448\ldots$. The masses
$p_j=2\rho^{N_j}$ form a probability distribution on the integers $N_j$.
Its raw moments are $m_k=\sum_j p_jN_j^k=-D^kF(\rho)$ for $k\ge1$.
Theorem~\ref{thm:lac} makes $\rho,m_1,m_2,\ldots$ algebraically independent.
For the displayed exponent sequence, the mean is $2.54434039308797664048\ldots$ and the variance is $0.90018704399788033852\ldots$, rounded at the last place. The variance $m_2-m_1^2$ is transcendental, and replacing $m_2$ by this variance preserves independence. More generally, the triangular changes to the mean and higher central moments or to the cumulants preserve independence in every finite order. There is also a counting interpretation of $1/\rho$. Allow two letters of each positive size $N_j$, and give a finite word the sum of its letter sizes. If $c_n$ counts the words of size $n$, summing over word lengths gives $\sum_{n\ge0}c_nz^n=1/F(z)$. For $|z|<\rho$ the generating function is analytic because $2\sum |z|^{N_j}<1$. It has a simple pole at $\rho$, since $F'(\rho)<0$. The Cauchy--Hadamard formula gives $\limsup c_n^{1/n}=1/\rho$, a transcendental growth rate (here $1/\rho=1.65331013984818334246\ldots$); the limsup also covers periodic sets of allowed sizes.

\section{Tribonacci reciprocal sums at geometric indices}
\label{sec:tribonacci}
Let $T_0=0$, $T_1=1$, $T_2=2$, and $T_{n+3}=T_{n+2}+T_{n+1}+T_n$.
Fix an integer $q\ge2$. For $a\in\overline{\mathbb Q}^{\times}$,
$\ell\in\mathbb Z_{\ge0}$, and $m\in\mathbb Z_{\ge1}$, define
\[
 S_q(a,\ell,m)=\sum_{k=0}^{\infty}\frac{k^\ell a^k}{T_{q^k}^{\,m}},
 \qquad 0^0=1 .
\]
\begin{theorem}\label{thm:tri}Every finite subfamily indexed by distinct triples $(a,\ell,m)$ is algebraically independent over $\K$.\end{theorem} The quantifier fixes $q$ throughout.
Tanaka's Theorem 3 concerns nongeometric index sequences. His Remark 2 identifies the geometric-index case with multiplicatively dependent characteristic roots as an open question in that work \cite{TanakaGeom}. The following argument addresses the particular Tribonacci recurrence.

\paragraph{Reduction to two variables.}
Write the roots of $X^3-X^2-X-1$ as $\alpha>1$, $\beta$, and $\gamma=\overline\beta$. Their product is $1$, and $|\beta|=|\gamma|=\alpha^{-1/2}$. The pair $(\beta,\gamma)$ is multiplicatively independent. Indeed, $\beta^u\gamma^v=1$ forces $u+v=0$ by absolute values. Moreover, $\beta/\gamma$ cannot be a root of unity. The irreducible cubic has a Galois conjugation carrying $\beta$ to $\alpha$ and $\gamma$ to a nonreal root, so a conjugate of that quotient has modulus greater than one.
The generating function $(z+z^2)/(1-z-z^2-z^3)$ has three simple poles. Consequently $T_n=A\alpha^n+B\beta^n+C\gamma^n$ with $ABC\ne0$. Set
\[
 U(x,y)=\frac{xy}{A+Bx^2y+Cxy^2},\qquad
 \Phi(x,y)=(x^q,y^q).
\]
Then $U(\beta^n,\gamma^n)=1/T_n$ for $n\ge1$. Define
\[
 J_{a,\ell,m}(x,y)=\sum_{k\ge0}k^\ell a^k U(\Phi^k(x,y))^m.
\]
These series converge near the origin. They satisfy
\[
 J_{a,\ell,m}=\mathbf 1_{\ell=0} U^m+
 a\sum_{j=0}^\ell \binom{\ell}{j} J_{a,j,m}\circ\Phi .
\]
For each $(a,m)$ this is a triangular block with diagonal $a$ and nonzero first subdiagonal. A finite family can be enlarged to include all the lower orders in its blocks.

\paragraph{The pole calculation.}
Nishioka's triangular-system lemma, in the form of Tanaka's Lemma 4 \cite{TanakaGeom}, says that a polynomial dependence among these functions forces a nonzero linear combination of the bottom entries of blocks with the same diagonal to be rational. Thus, for some $a\ne0$, a rational $V$ would satisfy
\[
 V-aV\circ\Phi=P(U),
\]
where $P$ is a nonzero polynomial with zero constant term. In the algebraic torus $(\mathbb C^\times)^2$, the pole divisor of $P(U)$ has the single irreducible component
$E=\{A+Bx^2y+Cxy^2=0\}$. Indeed, as a quadratic in $x$, its defining polynomial has discriminant $C^2y^4-4ABy$, with a simple zero at $y=0$. Its coefficients have no common factor.

The curves $\Phi^j(E)$ for $j\ge0$ are distinct. To see this, apply the surjective torus homomorphism $(x,y)\mapsto (x^2y,xy^2)$, which commutes with $\Phi$ and sends $E$ onto the line $A+Bu+Cv=0$. Under $(u,v)\mapsto(u^n,v^n)$ the image of this line has projective degree $n$. The restriction is generically one-to-one, since a coordinatewise pair of $n$th roots of unity preserving the line must be $(1,1)$. The inverse image of a generic projective line then meets the original line in $n$ points. Taking $n=q^j$ proves the assertion.

A pole of $V$ on any forward image $\Phi^j(E)$, $j\ge1$, would propagate indefinitely forward by the displayed equation, since its right side is regular at the generic point of each such curve. Thus $V$ has no pole on $\Phi(E)$ and must have a pole on $E$. Pulling back along any chain of irreducible preimages of $E$ now propagates poles indefinitely backward. Those curves are distinct because $E$ has distinct forward images. This contradicts finiteness of the pole divisor of $V$.

\paragraph{Specialization and a numerical instance.}
The multiplicative independence of $\beta,\gamma$ and their moduli below one give admissibility for $qI_2$. Along their forward orbit every denominator of $U$ is nonzero since $T_{q^k}>0$. Specializing the algebraically independent functions is therefore legitimate by the Mahler specialization theorem \cite[Cor.~3.5, Rem.~5.8 and Thm.~5.9]{AFfinite}. Equivalently, one may specialize far along the orbit and invert the triangular equations. The values at $(\beta,\gamma)$ are the claimed sums. For example,
\[
 \frac{S_2(1,0,2)}{S_2(1,0,1)^2}
     =0.4520046555845386649161706\ldots
\]
is transcendental. Algebraicity would give a polynomial relation between the two displayed sums.

\section{The three Mahler functions in Ide's question}
\label{sec:ide}
Let $d\ge2$ be an integer and put
\[
 F(z)=\sum_{k\ge0}z^{d^k},\quad
 G(z)=\prod_{k\ge0}(1-z^{d^k}),\quad
 H(z)=\sum_{k\ge0}\frac{z^{d^k}}{1-z^{d^k}} .
\]
Ide's introduction records the gap in the dependent-point literature for $G$ and $H$ as of 2023, together with earlier special cases \cite{IdeGeometric}. This question concerns the geometric exponents $d^k$. Ide's main results in the same paper treat nongeometric recurrences.

\begin{theorem}[Positive algebraic arguments]\label{thm:ide}
For a nonempty finite list of algebraic $0<a_i<1$, identify $i$ and $j$ whenever $a_i^{d^u}=a_j^{d^v}$ for some $u,v\ge0$. If $s$ is the number of resulting classes, then
\[
 \operatorname{trdeg}_{\overline{\mathbb Q}}
 \overline{\mathbb Q}\bigl(F(a_i),G(a_i),H(a_i):1\le i\le r\bigr)=3s.
\]
In particular all $3r$ entries are algebraically independent precisely when the forward orbits are pairwise disjoint.
\end{theorem}
The equations for one forward step are
\[
 F(z)-F(z^d)=z,\quad G(z)/G(z^d)=1-z,\quad
 H(z)-H(z^d)=\frac{z}{1-z}.
\]
These equations identify the fields of values at any two points in the same forward-orbit class.

\paragraph{Independence within a multiplicative dependence class.}
A finite collection of pairwise multiplicatively dependent positive arguments can be written $b^{e_i}$ with a common algebraic $0<b<1$ and $e_i\ge1$. This follows by clearing denominators in the rational ratios of their logarithms. Removing powers of $d$ from the $e_i$ by the equations leaves a set $M$ of distinct positive integers with $d\nmid m$ for each $m\in M$.

We claim the $3|M|$ functions $F(z^m),G(z^m),H(z^m)$ are algebraically independent over $\overline{\mathbb Q}(z)$. By the mixed additive and multiplicative criterion of Kubota (Tanaka's Lemma 5 in \cite{TanakaGeom}), it suffices to rule out rational solutions of
\begin{align*}
 V(z)-V(z^d)&=\sum_{m\in M}\left(c_m z^m+t_m\frac{z^m}{1-z^m}\right),\\
 B(z)/B(z^d)&=\prod_{m\in M}(1-z^m)^{n_m}
\end{align*}
for nonzero coefficient vectors, with $c_m,t_m\in\overline{\mathbb Q}$, $n_m\in\mathbb Z$ and $B\ne0$.

Here are details of both tests. For a root of unity $\xi$, let
$w(\xi)=\operatorname{Res}_{z=\xi}(V(z)\,dz/z)$. Pullback on residues in this normalization gives
\[
 w(\xi)-d^{-1}w(\xi^d)
 =-\sum_{m\in M}(t_m/m)\,\mathbf 1_{\xi^m=1}.
\]
The equation determines its finitely supported solution uniquely, since every nonzero root of unity has infinitely many iterated preimages. Thus $w$ depends only on the order of $\xi$, by invariance under automorphisms of the group of roots of unity. Every such finitely supported function is a finite linear combination of $e_j(\xi)=\mathbf 1_{\xi^j=1}$ by M\"obius inversion, and $e_j(\xi^d)=e_{dj}(\xi)$. These functions are linearly independent. In the chain $j,dj,d^2j,\ldots$ with $d\nmid j$, the right side can have a nonzero coefficient only at the first index. The last nonzero coefficient of $w$ on that chain would contribute a nonzero coefficient one step farther along. Thus all $t_m$ vanish. The polynomial part at infinity of $V$ then gives the same coefficient comparison in the monomials $z^j$, so all $c_m$ vanish. For the second test take $w(\xi)=\operatorname{ord}_{z=\xi}B(z)$. One obtains $w(\xi)-w(\xi^d)=\sum_m n_m e_m(\xi)$ and the same chain argument gives $n_m=0$ for every $m$.

\paragraph{Passage to values.}
The combined system is regular along the orbit of $b$, as all factors $1-b^{md^k}$ are nonzero. Nishioka's specialization theorem gives the same transcendence degree at $b$. Finally, the linear disjointness theorem of Adamczewski and Faverjon \cite[Theorem~1.1]{AFdisjoint} makes the degrees additive over the distinct multiplicative-dependence classes. This last step applies even when representatives from three or more classes satisfy a joint multiplicative relation. For $d=2$, the points $1/2$ and $1/8$ lie in one multiplicative class but have disjoint forward orbits, since $1$ and $3$ differ by no power of $2$. Their six values are algebraically independent.
In particular
\[
 \frac{G(1/2)}{G(1/8)}
  =0.40666198065076468037535465413485503216\ldots
\]
is transcendental. The proof uses the 2026 linear disjointness preprint as an input. The coverage in Ide's introduction is a historical statement about 2023.

\section{Joint derivatives of Stern's series and its twist}
\label{sec:sternjets}
Let $a_0=b_0=0$, $a_1=b_1=1$, and, for $n\ge1$, put
\[
\begin{aligned}
 a_{2n}&=a_n,& a_{2n+1}&=a_n+a_{n+1},\\
 b_{2n}&=-b_n,& b_{2n+1}&=-(b_n+b_{n+1}).
\end{aligned}
\]
Define $A(z)=\sum_{n\ge0}a_{n+1}z^n$ and
$B(z)=\sum_{n\ge0}b_{n+1}z^n$. Splitting the series into even and
odd powers yields, with $p(z)=1+z+z^2$,
\[
 A(z)=p(z)A(z^2),\qquad B(z)=2-p(z)B(z^2).
\]
Bundschuh proved the hypertranscendence of $A$, which means that every
finite list of its successive derivatives is algebraically independent
over $\mathbb C(z)$ \cite{sn-sternA}. Bundschuh and V\"a\"an\"anen proved
the same for $B$ and the algebraic independence of $A,B$. Immediately
after their Theorem~1.4 they asked about the two lists of derivatives
taken together \cite{sn-stern}. The following reconstruction uses the
ordinary difference Galois group of a finite system. These citations
document the historical question and the inputs to the argument.

\begin{theorem}[Joint Stern derivatives]\label{thm:sternjets}
For each $n\ge0$, the $2n+2$ functions
\[
 A,A',\ldots,A^{(n)},\quad B,B',\ldots,B^{(n)}
\]
are algebraically independent over $\mathbb C(z)$. At every algebraic
$\alpha$ with $0<|\alpha|<1$ their values are algebraically independent
over $\overline{\mathbb Q}$.
\end{theorem}

\begin{proof}
Set $\tau f(z)=f(z^4)$. Iterating the two equations gives a common
multiplier $u$:
\[
 \tau A=uA,\qquad \tau B=uB+v,\qquad
 u=\frac{1}{p(z)p(z^2)},\quad v=2\bigl(p(z)-1\bigr)u.
\]
Adjoin an indeterminate $L$ to the field of formal Puiseux Laurent
series $\bigcup_{s\ge1}\mathbb C((z^{1/s}))$. Extend $\tau$ by
$\tau L=4L$. Extend $\theta=z\,d/dz$ by $\theta L=1$ and put
$\delta=L\theta$. Then $\delta\tau=\tau\delta$. Work over the subfield
$K=(\bigcup_{s\ge1}\mathbb C(z^{1/s}))(L)$. Both the base and its
ambient field have $\tau$-constants $\mathbb C$. In fact, expansion
in $z^{1/s}$ over $\mathbb C(L)$ shows that an invariant nonzero
element has lowest exponent zero, since $\tau$ multiplies this exponent
by four. Its constant coefficient $r(L)$ satisfies $r(4L)=r(L)$.
A rational invariant of this scaling is constant, by its Laurent
expansion at $L=0$. Subtracting that constant leaves either zero or an
impossible invariant of positive lowest exponent.

For $R_n=\mathbb C[\varepsilon]/(\varepsilon^{n+1})$, write
\[
 Jf=\sum_{j=0}^n\frac{\delta^j f}{j!}\varepsilon^j,\qquad
 Y=\begin{pmatrix} JA&JB\\0&1\end{pmatrix}.
\]
The map $J$ preserves products and commutes with $\tau$. Consequently
\[
 \tau Y=\begin{pmatrix}Ju&Jv\\0&1\end{pmatrix}Y.
\]
Represent each entry by its matrix of multiplication in the basis
$1,\varepsilon,\ldots,\varepsilon^n$. The resulting matrix is a
fundamental solution of an ordinary linear difference system over $K$.
Its solution field $E=K(\delta^j A,\delta^j B:0\le j\le n)$ has no
new constants, by the ambient-field calculation. The standard
Picard--Vessiot dimension and restriction theorems therefore apply
\cite[\S1]{dsdhr}. Its group $G$ acts by $Y\mapsto YC$.
Since $C=Y^{-1}\sigma(Y)$ has constant coefficients and the
indicated affine matrix form, $G$ is a closed subgroup of
$R_n^\times\ltimes R_n$. Here $(c,t)$ denotes the map $x\mapsto cx+t$.

Projection to $R_n^\times$ is onto. Indeed the scalar subsystem with
solution $JA$ has transcendence degree $n+1$ over $K$. The extension
by $z^{1/s}$ is algebraic, $L$ is independent of the entire Puiseux
field, and the change from ordinary to $\delta$-derivatives is triangular
with leading coefficients $(Lz)^j$. Thus the known hypertranscendence
of $A$ gives the stated degree. Restriction onto this subsystem is
surjective. Its group is closed in the irreducible group $R_n^\times$
of dimension $n+1$ and has that same dimension. The same reasoning at
$\varepsilon=0$ uses the independence of $A,B$ to make the image of
$G$ the full two-dimensional affine group over $\mathbb C$.

Let $H$ be the kernel of the unit projection. This is a closed additive
subgroup of $R_n$, hence a complex vector subspace in characteristic
zero. It is stable under multiplication by every unit, by conjugation
and surjectivity of the unit projection. In the ordinary affine group,
the dilation $x\mapsto 2x$ and the translation $x\mapsto x+1$ have a
nonzero translation as commutator. Lift both to $G$. Their commutator
lies in $H$ because the unit group is abelian. Its translation
coefficient has nonzero constant term and is therefore a unit. The
unit action then puts every unit in $H$. The units span $R_n$ over
$\mathbb C$, so $H=R_n$. We conclude $\dim G=2(n+1)$. The
Picard--Vessiot dimension theorem gives the same value for
$\operatorname{trdeg}_K E$. Triangular conversion back to ordinary
derivatives proves the functional assertion.

For specialization use the ordinary derivative system over
$\overline{\mathbb Q}(z)$ together with the constant function $1$.
In the $j$th derivative of $\tau f$ the coefficient of
$f^{(j)}(z^4)$ is $(4z^3)^j$. The system and its inverse can therefore
have poles only at $0$ or zeros of $p(z)p(z^2)$. The latter all have
absolute value one. Thus the forward orbit of every indicated $\alpha$
is regular, and Nishioka's specialization theorem
\cite[Theorem 4.2.1]{sn-nish} gives the numerical assertion.
\end{proof}

In particular the Wronskian
\[
             W(\alpha)=A(\alpha)B'(\alpha)-A'(\alpha)B(\alpha)
\]
is transcendental at every such $\alpha$. At $\alpha=1/2$ this gives
\[
 W(1/2)=-5.14324842668610634571393163449568411372\ldots .
\]
If $W(\alpha)$ were algebraic, subtracting it from the displayed polynomial in four independent
values would give a nonzero algebraic relation. The same triangular
argument shows independence of $W(\alpha),\ldots,W^{(r)}(\alpha)$.
Indeed $W^{(j)}$ has leading term $A B^{(j+1)}$ in the successive
variables $B',\ldots,B^{(r+1)}$, and $A(\alpha)\ne0$ by the theorem.
 
\section{Weighted Lambert series, their relations, and a single crossing}
\label{sec:weightedlambert}
Fix an integer $d\ge2$ and a nonzero complex weight $a$. On $|z|<1$
put $h(z)=z/(1-z)$, $k(z)=z/(1+z)$, and define
\[
 X_{d,a}(z)=\sum_{t\ge0}a^t h(z^{d^t}),\qquad
 Y_{d,a}(z)=\sum_{t\ge0}a^t k(z^{d^t}).
\]
Double-exponential decay gives convergence locally uniformly. These are
the functions $G_d(a,z),F_d(a,z)$ of Bundschuh and V\"a\"an\"anen
\cite[equation (2.9)]{sn-lambert}. Two identities control their relations:
\begin{align}
 X_{d,a}(z)-aX_{d,a}(z^d)&=h(z),\label{eq:ls-shift}\\
 Y_{d,a}(z)&=X_{d,a}(z)-2X_{d,a}(z^2).\label{eq:ls-plus}
\end{align}
The first follows by shifting the summation index. The second follows
termwise from $k(z)=h(z)-2h(z^2)$.

\begin{theorem}[A fixed-base family with arbitrary weights]\label{thm:weightedlambert}
Write $\theta=z\,d/dz$. For a fixed $d$, the functions
\[
 \theta^j\!\left[X_{d,a}(z^u)\right],\qquad
       a\in\mathbb C^\times,\quad j\ge0,\quad u\ge1,\quad d\nmid u,
\]
are algebraically independent over $\mathbb C(z)$, in the sense of
every finite subfamily of distinct triples. At algebraic weights and
an algebraic point $\alpha$ with $0<|\alpha|<1$, the corresponding
values are algebraically independent over $\overline{\mathbb Q}$.
\end{theorem}

\begin{proof}
Write $U_{a,u,j}=\theta^j[X_{d,a}(z^u)]$. By
$\theta[f(z^d)]=d(\theta f)(z^d)$, equation~\eqref{eq:ls-shift} gives
\[
 U_{a,u,j}(z)-ad^jU_{a,u,j}(z^d)=\theta^j[h(z^u)].
\]
The constant-multiplier criterion of Nishioka
\cite[Theorems 3.2.1 and 3.2.2]{sn-nish}, applied as in
\cite[Section 3]{sn-bvmahler}, has the following consequence.
If a finite collection of these functions is algebraically dependent,
then in one block with common multiplier $\lambda=ad^j$ a nontrivial
constant linear combination is a rational function $g$. Equivalently,
\[
 g(z)-\lambda g(z^d)=H(z),\qquad
 H(z)=\sum_{ad^j=\lambda}c_{a,u,j}\,\theta^j[h(z^u)]
\]
with some coefficient nonzero.

Every pole of $H$ is a root of unity whose multiplicative order is
not divisible by $d$, because that order divides some $u$ with $d\nmid u$. Also, $H$
has at least one pole. Choose the largest $u$ with an active
coefficient and a primitive $u$th root $\zeta$. Terms indexed by
smaller exponents are analytic there. At this $u$ choose the largest
active $j$. The pole in $\theta^j[h(z^u)]$ has multiplicity $j+1$.
The multiplier determines $a=\lambda d^{-j}$ uniquely, so this
highest-order pole survives.

A pole of $g$ in $\mathbb C^\times$ away from roots of unity would
force all its iterated $d$th roots to be poles by the displayed
equation. Rationality excludes this infinite set. If $g$ has a
root-of-unity pole $\omega$ of maximal multiplicative order $N$, choose
a $d$th root $\eta$ of $\omega$ of exact order $dN$. To see the choice,
write $\omega=\exp(2\pi i b/N)$ with $(b,N)=1$ and choose $b'$ congruent
to $b$ modulo $N$ and coprime to $dN$ by avoiding the primes dividing
$d$ but not $N$. Then $\eta=\exp(2\pi i b'/(dN))$ works. At $\eta$,
both $g(z)$ and $H(z)$ are regular while $g(z^d)$ has a pole. This
contradiction excludes every pole of $g$ in $\mathbb C^\times$ and
leaves the established pole of $H$ unexplained. This proves functional
independence.

For algebraic weights the finite system, augmented by $1$, has
algebraic rational coefficients and nonzero constant diagonal
multipliers. All poles of its matrix and inverse in $\mathbb C^\times$
lie on the unit circle. Its power series coefficients belong to a
number field. Nishioka's specialization theorem
\cite[Theorem 4.2.1]{sn-nish} applies throughout the punctured open
disk at algebraic arguments.
\end{proof}

\subsection{Power substitutions and the dyadic exception}
For an integer $m>0$ let $c_d(m)$ be the unique $u$ for which
$m=d^t u$ with $t\ge0$ and $d\nmid u$. Iterating
\eqref{eq:ls-shift} gives
\[
 X_{d,a}(z^m)=a^{-t}\left(
 X_{d,a}(z^u)-\sum_{i=0}^{t-1}a^i h(z^{d^i u})\right).
\]
Thus all $X$-values with exponent $m$ reduce to its representative
$u$. Equation~\eqref{eq:ls-plus} expresses each $Y$-value through
the exponents $m$ and $2m$. Conversely, the pair $X_{d,a}(z^m)$,
$Y_{d,a}(z^m)$ recovers $X_{d,a}(z^{2m})$. All three assertions
continue to hold after applying any $\theta^j$.

In particular, let $S$ be a nonempty finite set of exponents, let
$\mathcal A$ be a finite set of distinct nonzero weights, and include
both types of series at every $m\in S$ and every $a\in\mathcal A$,
with derivatives through order $r$. The generated field has
transcendence degree
\[
  (r+1)|\mathcal A|\,
  \bigl|\{c_d(m):m\in S\cup\{2s:s\in S\}\}\bigr|.
\]
Indeed it equals the field generated by the indicated independent
representatives over $\mathbb C(z)$. The same degree holds over
$\overline{\mathbb Q}$ upon regular algebraic specialization with
algebraic weights. Ordinary jets and Euler jets give the same fields
there by triangular change of basis. The affine expressions above
also identify every polynomial relation: substitute the expressions
in independent representatives and test the resulting polynomial
identically. For $d=2$ the representatives reduce to the odd parts of
members of $S$.

Write $X_a=X_{2,a}$ and $Y_a=Y_{2,a}$. Equations
\eqref{eq:ls-shift} and \eqref{eq:ls-plus} give
\[
              (2-a)X_a+aY_a=2h.
\]
At $a=1$ this recovers the classical dyadic identity of
\cite[equation (1.6)]{sn-lambert}. At $a=2$ one has $Y_2=h$.
At $a=3$ one obtains $3Y_3=X_3+2h$. Therefore at $z=1/2$,
\[
 \frac{Y_3(1/2)}{X_3(1/2)}
   =\frac13+\frac{2}{3X_3(1/2)}
   =0.579597652270118\ldots .
\]
The positive series $X_3(1/2)$ is transcendental by
Theorem~\ref{thm:weightedlambert}. If the displayed ratio were
algebraic, solving its equation would make $X_3(1/2)$ algebraic,
since the ratio exceeds $1/3$. This proves the claimed transcendence
of the ratio.

\subsection{The unique positive crossing}
The equation $X_3(x)=2Y_3(x)$ on $(0,1)$ is equivalent to
$X_3(x)=4h(x)$. On this interval,
\[
 \frac{X_3(x)}{h(x)}
  =1+\sum_{t=1}^{\infty}
       \frac{3^t}{1+x^{-1}+\cdots+x^{-(2^t-1)}}.
\]
Each summand increases strictly with $x$. The quotient is continuous,
tends to $1$ as $x\downarrow0$ by its power series, and tends to
infinity as $x\uparrow1$. For the last assertion even its finite
partial sums have limits $1+\sum_{t=1}^M(3/2)^t$, unbounded as
$M$ increases. There is exactly one crossing $\xi$, numerically
$\xi=0.674745453983085\ldots$. If $\xi$ were algebraic, then
$X_3(\xi)=4\xi/(1-\xi)$ would be algebraic, contrary to
Theorem~\ref{thm:weightedlambert}. Thus the zero itself is
transcendental. The series identity determines the value at the
crossing, and strict monotonicity specifies which zero is meant.

\subsection{The weighted Fredholm series together with the Lambert series}
Write
\[
 F_{d,a}(z)=\sum_{k\ge0}a^kz^{d^k},\qquad F_{d,a}(z)-aF_{d,a}(z^d)=z .
\]
Here $F_{2,a}$ is the function $F_a$ in the dyadic family.
For fixed $d$, adjoin to the full independent family of
Theorem~\ref{thm:weightedlambert} all $F_{d,b}(z^u)$ with
$b\ne0$, $u\ge1$, $d\nmid u$, taking each pair $(b,u)$ once.
The enlarged family is algebraically independent over $\mathbb C(z)$.
For algebraic weights and any single complex algebraic $\alpha$ with
$0<|\alpha|<1$, its values remain algebraically independent. In
particular, for distinct nonzero algebraic $a_1,\ldots,a_s$, the
$2s$ numbers $F_{2,a_i}(\alpha),X_{2,a_i}(\alpha)$ are jointly
algebraically independent. Adding $Y_{2,a_i}(\alpha)$ gives precisely
the affine relations already displayed.

\begin{proof}
The same constant-multiplier criterion reduces a dependence to
$g(z)-\lambda g(z^d)=H(z)+\sum_u e_u z^u$ for a rational $g$, in
one common multiplier block. Here $H$ is the linear combination of
Lambert forcing terms in Theorem~\ref{thm:weightedlambert}.
If any coefficient from the Lambert series is nonzero, the
largest-exponent and highest-pole argument in that proof finds a pole
of $H$. The additional polynomial has no finite poles, so the same
backward propagation contradiction applies. Otherwise the right side
is a polynomial. The combination $g$ is analytic at zero, and a
nonzero finite pole would propagate to its infinitely many iterated
$d$th roots. Hence $g$ is a polynomial. If $g$ has positive degree
$N$, then $g(z)-\lambda g(z^d)$ has degree $dN$, whereas the largest
active exponent on the right is indivisible by $d$. A constant $g$
cannot supply any positive-degree term. Thus all $e_u$ must vanish.
For this block the Fredholm weight equals $\lambda$ and each $u$
occurs only once, so this contradicts nontriviality. The new equations
have polynomial forcing and nonzero constant multipliers, hence the
same specialization argument applies.
\end{proof}

There is an exact coincidence among the Fredholm derivatives:
\[
 \theta^j[F_{d,a}(z^u)]=u^j F_{d,ad^j}(z^u).
\]
Consequently a finite family containing $z$-derivatives of Fredholm
series must first merge the equal pairs $(ad^j,u)$. At one argument
$\alpha$, for weights in a finite set $\mathcal A$ and all Euler
derivatives through order $r$ of both $F$ and $X$, its transcendence
degree is
\[
  \bigl|\{ad^j:a\in\mathcal A,\ 0\le j\le r\}\bigr|
       +(r+1)|\mathcal A|.
\]
The formula also holds for the complete ordinary jets by their
triangular change of variables at $\alpha\ne0$.

For positive algebraic arguments the formula extends across the
forward-power classes. Within one multiplicative-dependence class
write the arguments as $\beta^m$ for a common algebraic
$0<\beta<1$ and positive integers $m$, by clearing the rational
denominators of their logarithmic ratios. Iterating the Fredholm
equation replaces $m=d^t u$ by $u$, up to a nonzero scalar and an
algebraic summand. The same holds for $X$ and its Euler derivatives
by \eqref{eq:ls-shift}. Distinct resulting exponents are independent
by the enlarged-family result. Across multiplicative-dependence
classes the value algebras are linearly disjoint by
\cite[Theorem~1.1]{AFdisjoint}, so the degrees add. In particular,
for the same weights at every point and both complete jets, the
displayed degree is multiplied by the number of distinct
forward-power classes. Two positive arguments $x,y$ represent the
same such class exactly when $x^{d^i}=y^{d^j}$ for some $i,j\ge0$.

At the unique real crossing $\xi$ described above the two common-scale
values are themselves transcendental:
\[
 X_3(\xi)=\frac{4\xi}{1-\xi}=8.298060239231769\ldots,\qquad
 Y_3(\xi)=\frac{2\xi}{1-\xi}=4.149030119615884\ldots .
\]
Indeed $\xi=X_3(\xi)/(4+X_3(\xi))$, so an algebraic value for
$X_3(\xi)$ would force the already transcendental $\xi$ to be
algebraic. The identity $Y_3(\xi)=X_3(\xi)/2$ gives the second
conclusion and shows the field generated by the three quantities has
transcendence degree one.

\section{The two limits of a continued fraction}
\label{sec:ds}

For \(0<|t|<1\), let \(E(t)\) and \(O(t)\) be the even- and odd-index limits, respectively, of
\([t,t^2,t^4,\ldots,t^{2^n}]\), where the last index is \(n\).
Dilcher and Stolarsky's series \(F,G\) have coefficients in \(\{0,1\}\), constant terms equal to one, and satisfy
\begin{equation}
 F(z)=G(z^2)+zF(z^4),\qquad G(z)=zF(z^2)+G(z^4).
 \label{eq:ds-rec}
\end{equation}
These recursions determine every positive-index coefficient from lower-index coefficients. The series have no zeros in the open unit disk, and the continued-fraction formulas are
\begin{equation}
 E(t)=t\frac{F(t^3)}{G(t^6)},\qquad
 O(t)=\frac{G(t^3)}{t^2F(t^6)}.
 \label{eq:ds-limits}
\end{equation}
These statements are proved in \cite[Sections 5--6]{ds09}. In 2014 Bundschuh and V\"a\"an\"anen asked about the algebraic independence of \(U(z)=F(z)/G(z^2)\) and \(V(z)=G(z)/F(z^2)\), and of their values at nonzero algebraic points in the disk \cite[Problem 1.15]{dsbv14}. Their subsequent rank-three theorem for the underlying series \cite[Theorems 1.2--1.3]{dsbv15} is the arithmetic starting point here.

\begin{theorem}[Rank at powers of one algebraic point]
\label{thm:ds-rank}
Let \(\alpha\) be algebraic with \(0<|\alpha|<1\). For a nonempty finite set \(M\) of positive integers let
\[
 U_M=\{m/2^{v_2(m)}:m\in M\},\qquad h=|U_M|.
\]
For each integer \(N\geq0\), the field over \(\overline{\mathbb Q}\) generated by
\[
 E^{(j)}(\alpha^m),\quad O^{(j)}(\alpha^m)
 \qquad(m\in M,\ 0\leq j\leq N)
\]
has transcendence degree \(2h(N+1)\). The field of the corresponding functions of a variable \(t\) has the same transcendence degree over \(\mathbb C(t)\).
\end{theorem}

\paragraph{A fundamental matrix.}
Use a formal cube root and put
\[
 B(z)=\begin{pmatrix}z+z^2&1\\z^2&1\end{pmatrix},\qquad
 Y(z)=\begin{pmatrix} F(z)&z^{-1/3}G(z)\\G(z^2)&z^{2/3}F(z^2)\end{pmatrix}.
\]
Equation \eqref{eq:ds-rec} gives \(Y(z)=B(z)Y(z^4)\) by substitution in each column. Also \(\det B(z)=z\). Set \(D(z)=G(z)G(z^2)-zF(z)F(z^2)\). Equation \eqref{eq:ds-rec} gives \(D(z)=F(z)G(z^4)-zF(z^4)G(z)=D(z^2)\). Because \(D(0)=1\), analyticity at zero gives
\begin{equation}
 G(z)G(z^2)-zF(z)F(z^2)=1,\qquad \det Y(z)=-z^{-1/3}.
 \label{eq:ds-det}
\end{equation}
This also proves the determinant identity of \cite[Proposition 5.2]{ds09} directly.

Let \(K=\bigcup_{e\geq1}\mathbb C(z^{1/e})\), with automorphism \(\sigma(z^{1/e})=z^{4/e}\). Its field of fixed elements is \(\mathbb C\), as is seen by comparing the orders of the nonconstant terms in a Puiseux expansion. The same comparison works for the ambient field of formal Puiseux series. Consequently the field generated by a fundamental matrix in this ambient field is an ordinary Picard--Vessiot field. Its transcendence degree over \(K\) equals the dimension of its Galois group. By \eqref{eq:ds-rec}, the entries of \(Y\) generate the same field over \(K\) as \(F(z),F(z^4),G(z),G(z^4)\). The latter field has transcendence degree three by \cite[Theorem 1.2]{dsbv15} and the algebraicity of \(K/\mathbb C(z)\). The determinant lies in \(K\), so the Galois group is a dimension-three closed subgroup of the connected, three-dimensional group \(\mathrm{SL}_2\). It equals \(\mathrm{SL}_2\). Substitution \(z\mapsto z^a\) is an automorphism of the Puiseux field carrying \(K\) onto itself and commuting with \(\sigma\). The same conclusion holds for \(Y(z^a)\) for every positive integer \(a\).

\paragraph{A short proof for the original pair.}
For any matrix \(Y=\left(\begin{smallmatrix}p&q\\c&d\end{smallmatrix}\right)\) above, write \(r=p/c\), \(s=q/d\). The entries are rational in \(r,s,c\) and the known determinant, since \(d=\det Y/[c(r-s)]\). Hence the rank of the four entries is at most one more than that of the two slopes. Bundschuh and V\"a\"an\"anen's function and value theorems give rank three in both settings, including after extraction of a cube root of an algebraic argument. Their results therefore imply rank two for \(r,s\). As \(r=U(z)\) and \(s=z^{-1}V(z)\), this proves both assertions for \(U,V\) in the 2014 question. The following argument handles powers with different odd parts and arbitrary orders of differentiation.

\paragraph{Independence for distinct odd exponents.}
We claim that for distinct positive odd \(a_1,\ldots,a_h\), the block system with matrices \(B(z^{a_i})\) has ordinary group \(\mathrm{SL}_2^h\).
Its group \(H\) projects onto each factor by restriction to the corresponding Picard--Vessiot subfield. The Lie algebra \(\mathfrak h\) projects onto each copy of \(\mathfrak{sl}_2\). Its solvable radical is zero, since each image of the radical is a solvable ideal of \(\mathfrak{sl}_2\). Each coordinate projection is therefore an isomorphism on one simple summand and zero on the others. Indeed the image of each summand is an ideal, and images of distinct summands commute. If \(\mathfrak h\) is smaller than the full product, two coordinates come from the same summand. Their projection is the graph of a Lie algebra isomorphism. Every element of \(H\) normalizes this graph, so the adjoint representations on these two coordinates are intertwined by this isomorphism. Picard--Vessiot descent then gives a rational isomorphism over \(K\) between the two adjoint difference systems. Explicitly, conjugating the constant intertwiner on both sides by their two adjoint fundamental matrices gives a matrix fixed by \(H\), hence with entries in \(K\).

Here is the obstruction to that isomorphism. Write the two exponents as \(a,b\), ordered so that \(a\nmid b\). Choose \(e\) so that a putative rational gauge \(P\) and its inverse are defined over \(\mathbb C(w)\), \(z=w^e\). The coefficient matrices \(\operatorname{Ad}B(w^{ea})\), \(\operatorname{Ad}B(w^{eb})\), and their inverses are regular on \(\mathbb C^\times\). Here \(\operatorname{Ad}A\) acts on traceless matrices by \(X\mapsto AXA^{-1}\). The gauge identity expresses \(P(w^4)\) as \(P(w)\) multiplied on each side by matrices regular and invertible at every nonzero \(w\). Such multipliers preserve regularity in both directions. The pole set of \(P\) in \(\mathbb C^\times\) thus contains all fourth roots of each of its elements. A finite nonempty set cannot have this property. Applying the same argument to \(P^{-1}\) shows that \(P\) is everywhere invertible on \(\mathbb C^\times\).

Write \(e=2^s e_0\) with \(e_0\) odd, and take \(w\) of exact order \(e_0a\). Then \(w^{ea}=1\) and \(\beta=w^{eb}\ne1\). Choose \(L\geq1\) with \(w^{4^L}=w\). Iterating the gauge along this cycle would make the adjoints of the two matrices
\[
 C^L,\quad M_\beta=B(\beta)B(\beta^4)\cdots B(\beta^{4^{L-1}}),
 \qquad C=B(1)=\begin{pmatrix}2&1\\1&1\end{pmatrix},
\]
conjugate. For every nontrivial odd-order root \(\omega\), the entrywise modulus \(|B(\omega)|\) is bounded by \(C\), strictly in the upper-left entry since \(|1+\omega|<2\). This applies to every point of the \(\beta\)-cycle. Consequently
\( |M_\beta|\leq\prod_{j=0}^{L-1}|B(\beta^{4^j})|\leq C^L\),
with a strict upper-left inequality in the second comparison. Since \(C^L\) has positive entries, strict Perron--Frobenius comparison gives
\(\rho(M_\beta)<\rho(C)^L\), where \(\rho\) denotes spectral radius. Both determinants have modulus one. For a two-by-two matrix \(A\) with this property, \(\rho(\operatorname{Ad}A)=\rho(A)^2\), as its adjoint eigenvalues are \(1,\lambda_1/\lambda_2,\lambda_2/\lambda_1\). The unequal adjoint spectral radii rule out the gauge, proving the claim.

\paragraph{Derivatives and specialization.}
We give the differential step explicitly because independence of the values alone does not give independence of their derivatives. Work over the parameterized Picard--Vessiot base of \cite{dsdhr}, adjoining \(\log z\) and differentially closed constants \(\widetilde{\mathbb C}\), with commuting derivation \(\delta=(\log z)z\,d/dz\) on the original functions. The displayed fundamental matrices embed in that extension by taking them in \(\bigcup_{e\geq1}\widetilde{\mathbb C}(\ell)((z^{1/e}))\), where \(\ell=\log z\), \(\sigma(\ell)=4\ell\), and \(\delta(\ell)=\ell\). This field has fixed field \(\widetilde{\mathbb C}\). Indeed, a fixed series can have only exponent zero, by comparing exponents under multiplication by four, and a rational function of \(\ell\) invariant under \(\ell\mapsto4\ell\) is constant. Thus the differential field generated by these matrices over the parameterized base introduces no new difference constants and is a parameterized Picard--Vessiot field. For each block, the coefficient of \(\sigma(Y)=A Y\) is \(A=B(z^{a_i})^{-1}\), with determinant \(z^{-a_i}\). The ordinary group contains \(\mathrm{SL}_2\), so \cite[Theorem 4.1]{dsdhr} places \(\mathrm{SL}_2(\widetilde{\mathbb C})\) in the parameterized group of that block. Its determinant is already fixed by \eqref{eq:ds-det}, giving equality. For the full system the parameterized group is Zariski dense in the ordinary product (adjoining the logarithm can remove only a multiplicative quotient by \cite[Proposition 1.6]{dsdhr}, whereas this product is perfect, and extending the constants extends the group by scalars as in \cite[Section 2]{dsdhr}). Its projections are the individual parameterized groups by restriction. For \(h\geq2\), \cite[Theorem 4.5]{dsdhr} identifies the differential closure of the derived connected group with the product of the corresponding groups of the factors. These factors are the full copies of \(\mathrm{SL}_2(\widetilde{\mathbb C})\), which are connected and perfect. Thus the parameterized group is the full product, of differential dimension \(3h\). For \(h=1\) the same dimension statement follows from the single-block result.

For a block \(Y=\left(\begin{smallmatrix}p&q\\ c&d\end{smallmatrix}\right)\), put \(r=p/c\) and \(s=q/d\). Then
\begin{equation}
  d=\frac{\det Y}{c(r-s)},\qquad p=rc,\qquad q=sd.
  \label{eq:ds-coordinates}
\end{equation}
Thus \((r,s,c)\) for the \(h\) blocks generate the differential field of their entries over the parameterized base. By the differential dimension theorem for a parameterized Picard--Vessiot extension, its differential transcendence degree is \(3h\). A differential generating set of exactly this size must be differentially algebraically independent. In particular all finite lists of derivatives of these coordinates are algebraically independent. Switching between \(\delta\), \(\theta=z\,d/dz\), and ordinary derivatives at a fixed finite order uses invertible triangular transformations over the base extended by \(\log z\), so the same independence holds for ordinary derivatives over \(\mathbb C(z)\).

To obtain arithmetic statements take the blocks \(Y(t^{3u})\), \(u\in U_M\). The exponents \(3u\) are distinct odd integers, so the preceding argument applies in the variable \(t\). For each block, multiply its entries by \(t^u\) to clear the only possible pole at zero. Equation \eqref{eq:ds-rec} gives rational coefficients for these Taylor series. The coefficient matrix of the gauged Mahler system and its inverse have no pole on \(\mathbb C^\times\). Adjoin the \(\theta\)-derivatives up to order \(N\) of every entry. Differentiating the system, using \(\theta\sigma=4\sigma\theta\), gives a triangular block system with invertible diagonal blocks. Its coefficient matrix and inverse still have no pole on \(\mathbb C^\times\). The orbit \(\{\alpha^{4^k}:k\geq0\}\) therefore consists of regular points. Nishioka's specialization theorem, in the form of \cite[Lemma 2.1]{dsbv15}, bounds the transcendence degree of the specialized entry field below by its functional transcendence degree \(3h(N+1)\). Multiplication by \(t^u\) and switching to ordinary derivatives preserve the fields of finite jets over the base, both before and after specialization.

The nonvanishing of \(F,G\) and of the determinant guarantees that \(c,d,r-s\) are nonzero at \(\alpha\). Formula \eqref{eq:ds-coordinates}, differentiated through order \(N\), shows that the field of specialized entries is generated by the \(3h(N+1)\) values from \((r,s,c)\) and their derivatives. They are therefore algebraically independent. For the block \(Y(t^{3u})\) its two slopes multiplied by \(t^u\) give \(E(t^u),O(t^u)\) by \eqref{eq:ds-limits}. Multiplication by \(t^u\) and the chain rule for \(t\mapsto t^u\) induce invertible triangular changes of finite jets at \(\alpha\ne0\). This proves the rank formula for odd exponents.

For the remaining exponents, removing the initial partial quotient gives
\[
 E(t)=t+1/O(t^2),\qquad O(t)=t+1/E(t^2).
\]
These are also identities of meromorphic functions by \eqref{eq:ds-rec}--\eqref{eq:ds-limits}. The formulas and their rational inverses identify the fields of each pair of finite jets at \(t\) and \(t^2\) over \(\overline{\mathbb Q}\) when \(t\) is algebraic, \(0<|t|<1\). All divisions are legitimate by \eqref{eq:ds-limits} and the nonvanishing just proved. The same identities hold over the function field. Iteration replaces each \(m\) by its odd part and completes the proof of Theorem~\ref{thm:ds-rank}. \hfill\(\square\)

\paragraph{A single real zero.}
From \eqref{eq:ds-det} and \eqref{eq:ds-limits},
\[
 \Delta(t):=O(t)-E(t)=\frac{1}{t^2 F(t^6)G(t^6)}.
\]
There is exactly one \(\xi\in(0,1)\) satisfying
\begin{equation}
 \Delta(\xi)\Delta(\xi^3)=1.\label{eq:ds-zero}
\end{equation}
Indeed, on this interval the denominator of \(\Delta\) is positive and strictly increasing. Each series has infinitely many nonzero coefficients by \eqref{eq:ds-rec} (for instance \(F\) propagates exponents via \(n\mapsto4n+1\), and \(G\) via \(n\mapsto4n\) starting at \(1\)). Thus \(\Delta\) decreases continuously from \(+\infty\) to zero. The product in \eqref{eq:ds-zero} has the same monotonicity and endpoint limits. If \(\xi\) were algebraic, the case \(M=\{1,3\},N=0\) of Theorem~\ref{thm:ds-rank} would make the four numbers \(E(\xi),O(\xi),E(\xi^3),O(\xi^3)\) algebraically independent. Equation \eqref{eq:ds-zero} gives the nonzero polynomial relation \((X_2-X_1)(X_4-X_3)-1=0\) on them. Hence \(\xi\) is transcendental.

Numerically \(0.8701938455851449618<\xi<0.8701938455851449619\).
The enclosure can be reproduced with rational arithmetic. At either rational endpoint \(t\), truncate each series through degree \(80\), using \eqref{eq:ds-rec} for its coefficients. For \(x=t^6,t^{18}\) the corresponding tail is bounded by \(x^{81}/(1-x)\). Substituting these upper and lower bounds in \(1/[t^8 F(t^6)G(t^6)F(t^{18})G(t^{18})]-1\) gives a positive value at the lower endpoint and a negative value at the upper endpoint.

\section{Chang's bilateral Fibonacci sums}\label{sec:chang}
Let $F_0=0,F_1=1,F_{n+2}=F_{n+1}+F_n$ and extend to negative
indices by $F_{-n}=(-1)^{n+1}F_n$. For the rest of this section put
\[
 \phi=(1+\sqrt5)/2,\qquad s=-\sqrt[3]{\phi},\qquad
 x=s/\phi=-\phi^{-2/3},\qquad T_\ell=\sum_{n\ne0}(s^n/F_n)^\ell.
\]
The cube root of $\phi$ is positive. The sums converge absolutely for each integer $\ell\ge1$. Chang's Theorem 1.1 gives $T_1^2=T_2$ and records $6.1469$ as an approximation to the common value \cite{Chang}. His concluding question asks for algebraic weights $\psi_\ell$ producing additional dependences among $\sum_{n\ne0}\psi_\ell^n/F_n^\ell$ for even $\ell$. The broad existential interpretation already admits $\psi_\ell=1$: the even bilateral sums are twice their one-sided counterparts, and the 2007 theorem of Elsner, Shimomura and Shiokawa puts them in a field of transcendence degree three \cite{ESS}. Our formulas concern the fixed choice $\psi_\ell=s^\ell$, which also allows odd $\ell$. We make no claim of bibliographic priority for the classification.

\subsection{A polynomial certificate for the common value}
The ratio of the two roots of the Fibonacci characteristic polynomial, smaller over larger, is $-\phi^{-2}=x^3$. Hence $x$ is the unique root in $(-1,0)$ of $x^6+3x^3+1=0$. Binet's formula gives for $n>0$
\[
 \frac{s^n}{F_n}=\frac{\sqrt5\,x^n}{1-x^{3n}},
 \qquad \frac{s^{-n}}{F_{-n}}=-\frac{\sqrt5\,x^{2n}}{1-x^{3n}} .
\]
Define the cubic theta function
\[
 A(q)=\sum_{m,n\in\Z}q^{m^2+mn+n^2}
     =1+6\sum_{n\ge1}\frac{q^n-q^{2n}}{1-q^{3n}}.
\]
The second equality is the Borwein--Borwein--Garvan identity
\cite[Eq.~(2.25)]{BBG}. It gives
\begin{equation}\label{eq:ChangC}
 T_1=\frac{\sqrt5}{6}(A(x)-1),\qquad
 C=T_2=\frac5{36}(A(x)-1)^2.
\end{equation}
The identity between the two sums also follows directly at the level of functions. In the notation of \eqref{eq:Hl}, expansion of the geometric square gives
\[
 u_2(d)=\frac{d\,\mathbf 1_{3\nmid d}-\chi(d)}{3},\qquad
 H_2=\frac{3E_2(q^3)-E_2(q)+2-4A}{72},
\]
where $\chi$ is the nontrivial character modulo three. The modular form $3E_2(q^3)-E_2(q)$ has weight two and constant term two. The weight-two space for $\Gamma_1(3)$ is spanned by $A^2$ by the ring description below. Thus $3E_2(q^3)-E_2(q)=2A^2$ and $H_2=(A-1)^2/36=H_1^2$. Binet's formula now gives $T_2=T_1^2$.

Write $B(q)=b(q)^3$, where the cubic theta $b(q)=\sum_{m,n}\omega^{m-n}q^{m^2+mn+n^2}$ uses a primitive cube root of unity $\omega$.
Matsuda's identities \cite[Thm.~5.1]{Matsuda}, with normalized Eisenstein series, read
\[
 U=E_4=9A^4-8AB,\qquad V=E_6=-27A^6+36A^3B-8B^2.
\]
Multiplication and subtraction eliminate $B$ without division:
\[
 27 A^8-18UA^4-8VA^2-U^2=0 .
\]
If $C$ were algebraic, \eqref{eq:ChangC} would give $A(x)=c\in\K$. The polynomial $27c^8-18c^4X-8c^2Y-X^2$ has coefficient $-1$ at $X^2$, yet it would vanish at $(E_4(x),E_6(x))$. Nesterenko's Corollary 2 \cite{Nest} gives algebraic independence of $E_2(x),E_4(x),E_6(x)$ because $0<|x|<1$ and $x$ is algebraic. This proves transcendence of both $T_1$ and $C$. In fact the ideal of all relations on $(T_1,T_2)$ in $\K[X,Y]$ is $(Y-X^2)$, by division in $Y$ and the transcendence of $T_1$.

\subsection{The specialization map}
Write $q=e^{2\pi i\tau}$. We use the ring of holomorphic modular forms for $\Gamma_1(3)$, allowing finite sums of different weights. Over $\C$ it is a polynomial ring on generators of weights $1$ and $3$, with algebraic Fourier coefficients. Adjoining $E_2(q)$ gives the polynomial ring of quasimodular forms \cite[\S4]{Clery}. For example one can use $A$ and $\eta(3\tau)^9/\eta(\tau)^3$ as the modular generators. At any fixed weight the distinct monomials are a basis. A finite coefficient matrix of this basis has full column rank, since otherwise the series would be dependent. Thus a form of that weight with algebraic Fourier coefficients has algebraic coordinates in the basis. We use the resulting $\K$-algebra, denoted $\mathcal M$, and $\mathcal M[E_2]$. The Eisenstein expansions below have algebraic constants in each weight. The same rings result upon adjoining $E_2(q^3)$ in place of $E_2(q)$ since $E_2(q)-3E_2(q^3)$ is modular of weight two.
\begin{lemma}\label{lem:evaluation}
At every algebraic $q_0$ with $0<|q_0|<1$, evaluation is injective on $\mathcal M[E_2]$. In particular it preserves all algebraic relations and all algebraic independences over $\K$ among its elements.
\end{lemma}
\begin{proof}
The field of values of the three generators contains $E_2(q_0), E_4(q_0), E_6(q_0)$, three algebraically independent elements by \cite{Nest}. The three generator values must therefore be algebraically independent. Evaluation on their polynomial ring is injective, and a relation on evaluated elements is equivalent to the evaluation of the corresponding polynomial in those elements.
\end{proof}

\subsection{Lambert kernels at all powers}
Define functions for $|q|<1$ by
\begin{align}
 H_\ell(q)
 &=\sum_{n\ge1}\frac{q^{\ell n}+(-1)^\ell q^{2\ell n}}{(1-q^{3n})^\ell}
 =\sum_{d\ge1}u_\ell(d)\frac{q^d}{1-q^d},\label{eq:Hl}\\
 u_\ell(d)&=[y^d]\frac{y^\ell+(-1)^\ell y^{2\ell}}{(1-y^3)^\ell},
 \qquad T_\ell=5^{\ell/2}H_\ell(x).\notag
\end{align}
In particular $H_\ell(q)=q^\ell+O(q^{\ell+1})$. Put $p_a(t)=\binom{(t-a)/3+\ell-1}{\ell-1}$ for $a=\ell,2\ell$, interpreting the binomial as a polynomial. On the congruence class $d\equiv a\pmod3$ its polynomial value gives the required positive-index coefficients, including the zero coefficients for $0<d<a$. Thus $u_\ell$ is the sum of $p_\ell$ and $(-1)^\ell p_{2\ell}$ on their respective supports. The reflection identity $p_\ell(-d)=(-1)^{\ell-1}p_{2\ell}(d)$ makes this extension odd. In particular the coefficient of $d^j$ has parity $(-1)^{j+1}$ under $d\mapsto-d$ in the residue group.

Let $\chi$ be the nontrivial Dirichlet character modulo 3. For $3\nmid\ell$, each even power $d^{2i}$ thus occurs with a constant multiple of $\chi(d)$, and each odd power $d^{2i+1}$ with a constant multiple of $\mathbf1_{3\nmid d}$. Their Lambert series are constants plus modular forms of weights $2i+1$ and $2i+2$, by the divisor-sum Fourier expansions of Eisenstein series \cite[Prop.~2.2 and Rem.~2.4]{BG}. In weight one the series is proportional, up to its constant term, to $A$. For even weight $w$ the principal-character term gives $E_w(q)-3^{w-1}E_w(q^3)$ up to scaling and a constant. At $w=2$ this very combination is modular. These statements follow as well by subtracting the terms with $3\mid d$ from the full Lambert series.

If $\ell=3k$, only $d=3t$ occurs. Set
\[
 P_k(t)=\binom{t+2k-1}{3k-1},\qquad u_{3k}(3t)=P_k(t)-P_k(-t).
\]
The product for $P_k$ has consecutive shifts $-(k-1),\ldots,2k-1$, with a single zero shift. The coefficient of $t$ in the difference is
\begin{equation}\label{eq:linear-anomaly}
 a_k=2(-1)^{k-1}\frac{(k-1)!(2k-1)!}{(3k-1)!}.
\end{equation}
Since $\sum_{t\ge1}tq^{3t}/(1-q^{3t})=(1-E_2(q^3))/24$, this proves
\begin{equation}\label{eq:split}
 H_\ell=M_\ell+c_\ell E_2(q^3),\qquad M_\ell\in\mathcal M,\qquad
 c_\ell=\begin{cases}0&3\nmid\ell,\\-a_k/24&\ell=3k.\end{cases}
\end{equation}
For example $H_3=(1-E_2(q^3))/24$. When $\ell=6$ one obtains, by expanding the binomial,
\[
 u_6(d)=\mathbf1_{3\mid d}(d^5/14580+d^3/324-d/30).
\]
The last term contributes $(E_2(q^3)-1)/240$, or $E_2(q)/720$ plus a modular term and a constant. This exhibits the third generator in the first even case.

\subsection{Distinguishing highest weights}
We give the comparison in detail because counting the generators of a ring supplies only an upper bound for the rank of a chosen subfamily. If $F,G$ are positive-weight homogeneous forms of weights $u,v$ in the polynomial ring $\mathcal M$, algebraic dependence implies that $F^v/G^u$ is constant. Indeed, separate any polynomial relation by total weight. In one nontrivial weighted component, division by one monomial yields a univariate relation on $F^{v/g}/G^{u/g}$, where $g=\gcd(u,v)$. The coefficient field is algebraically closed. Also, independence of highest-weight components implies independence of the original mixed-weight expressions: the highest weighted piece of a putative relation would survive upon substitution.

By Chang's identity we may represent index 1 by 2 for computing transcendence degrees. With $G_w=E_{w,1,\chi}$ denoting the character Eisenstein series whose nonconstant coefficients are proportional to $\sum_{d\mid n}\chi(d)d^{w-1}$, the highest-weight components of the $M_\ell$, up to nonzero constants, are
\begin{center}
\begin{tabular}{@{}ll@{}}\toprule
Index & Highest-weight component\\\midrule
$3\nmid\ell$, even & $D_\ell=E_\ell(q)-3^{\ell-1}E_\ell(q^3)$\\
$3\nmid\ell$, odd, $\ell\ge5$ & $G_\ell$\\
$3\mid\ell$, $\ell\ge6$ & $L_w=E_w(q^3)$, $w=2\lfloor\ell/2\rfloor$\\\bottomrule
\end{tabular}
\end{center}
For the first two rows the leading terms of the residue polynomials occur on distinct residue classes. In the last row $P_k(t)-P_k(-t)$ has degree $3k-1$ for even $k$. For odd $k>1$ the top even term disappears and the next term survives because the sum of the shifts is $k(3k-1)/2$. The resulting modular weights for the successive multiples $6,9,12,15,\ldots$ are $6,8,12,14,\ldots$, strictly increasing. At $\ell=3$, $M_3$ is a constant.

Here are the three comparisons needed for the table. For a form $f$ with nonzero constant term and weight $w$, consider the coefficient of its first nonconstant term in $f/f(0)$ divided by $w$, with $f(0)$ denoting its coefficient at $q^0$. Taking a positive integral power leaves this normalized quantity unchanged when the weight is adjusted. For $D_w$ the first exponent is 1, and the quantity is $-2/[B_w(1-3^{w-1})]$. For $L_w$ the first exponent is 3, and the quantity is $-2/B_w$. Here $B_w$ is the Bernoulli number. The numbers $|B_w|$ increase strictly through even $w\ge6$, and $(3^{w-1}-1)|B_w|$ increases strictly through even $w\ge2$. Indeed
\[
 \frac{|B_{w+2}|}{|B_w|}
 =\frac{(w+1)(w+2)}{4\pi^2}\frac{\zeta(w+2)}{\zeta(w)}
 >\frac{(w+1)(w+2)}{4\pi^2\zeta(w)}.
\]
At $w\ge6$ the lower bound exceeds one, since $\zeta(6)<6/5$.
For $w\ge2$ the additional factor $(3^{w+1}-1)/(3^{w-1}-1)$ exceeds nine, and $9\cdot12/(4\pi^2\zeta(2))>1$. Thus distinct weights in either row cannot give proportional equal-weight powers, and the first exponents distinguish the two rows from each other.

For the odd-weight row the Fricke involution $W_3$ takes $G_w$ to a nonzero multiple of $E_{w,\chi,1}$ \cite[Cor.~4.4]{EKS}. The latter has zero constant term and nonzero $q$ coefficient. Hence every $G_w$ has order one at the cusp $0$, whereas its order divided by its weight is $1/w$. On the other hand both $D_w$ and $L_w$ have nonzero constant term at that cusp. In fact $L_w|W_3=3^{-w/2}E_w(q)$, and $D_w|W_3=3^{w/2}E_w(q^3)-3^{w/2-1}E_w(q)$. For $w=2$ in the latter formula the anomalous terms cancel. This separates every odd-row form from the other rows and separates distinct weights within the odd row. Every two distinct entries that occur after representing 1 by 2 therefore have independent top forms, except that index 3, already isolated above, is handled directly by \eqref{eq:split}.

\subsection{The exact rank and its arithmetic consequences}
\begin{theorem}\label{thm:chang-rank}
For a finite set $S$ of positive integers put $m(S)=|S|-\mathbf1_{\{1,2\}\subseteq S}$. Then
\[
 \trdeg_{\K}\K(T_\ell:\ell\in S)
 =\min\{m(S),\ 2+\mathbf1_{S\cap3\Z\ne\varnothing}\}.
\]
\end{theorem}
\begin{proof}
Replacing index 1 by 2, even when only 1 was present, preserves the transcendence degree by $T_2=T_1^2$. The resulting set has $m(S)$ elements. The algebraic scalar factors in \eqref{eq:Hl} have no effect. For an index set disjoint from $3\Z$, all $H_\ell$ are modular. Any selected pair has independent top forms, and every singleton is nonconstant. The modular ring has two generators.

Otherwise take the smallest multiple $j$ of 3 in the set. For every other index replace $H_\ell$ by $H_\ell-(c_\ell/c_j)H_j$. These new functions are modular and together with $H_j$ generate the original field. The highest component of each changed function is preserved, since a nontrivial subtraction occurs only at a larger multiple of 3; the increasing weights were determined above. Each unmodified modular function retains its own highest component. Every pair of these modular functions is independent, and each single function is nonconstant. Finally, $H_j$ is transcendental over the entire modular function field of fractions $\operatorname{Frac}(\mathcal M)$, by \eqref{eq:split} and polynomial independence of the quasimodular generator. This adds exactly one to the modular rank. Lemma~\ref{lem:evaluation} transfers the result to the common argument $x$.
\end{proof}

In particular all individual $T_\ell$ are transcendental. For the even subfamily the rank is $\min(|S|,2+\mathbf1_{S\cap6\Z\ne\varnothing})$. The complete positive even family has rank three and its subfamily with $3\nmid\ell$ has rank two. The same conclusions for infinite sets follow by taking suprema over finite subsets. Examples of independent triples are $(2,4,6)$, $(3,5,7)$, $(3,9,15)$ and $(6,12,18)$, where a triple denotes its three indices. The triples $(2,4,8)$ and $(5,7,11)$ are dependent with every pair independent. The triple $(1,2,3)$ has rank two.
\begin{corollary}
The family $1,T_1,T_2,\ldots$ is linearly independent over $\K$. Every nontrivial finite algebraic linear combination of the $T_\ell$ is transcendental.
\end{corollary}
\begin{proof}
The distinct positive orders $\ell$ of the series $H_\ell$ in \eqref{eq:Hl} isolate the smallest selected index after testing the constant term. This proves independence of $1$ and the functions. Apply Lemma~\ref{lem:evaluation} and multiply each coordinate by its nonzero algebraic scale.
\end{proof}
For a nonlinear example, independence of $(T_6,T_{12},T_{18})$ makes the moment determinant $T_6T_{18}-T_{12}^2$ transcendental. If it were $a\in\K$, the polynomial $XZ-Y^2-a$ would contradict Theorem~\ref{thm:chang-rank}.

\appendix
\section{Numerical checks of the Fibonacci formulas}\label{sec:numerics}
The numbers in the table use \eqref{eq:Hl} at the negative real argument $x=-\phi^{-2/3}$. To bound the error after the first $N$ summands, put $r=|x|$. For every $n>N$, the two terms in the numerator have total modulus at most $2r^{\ell n}$ and the denominator has modulus at least $(1-r^{3(N+1)})^\ell$. Summing the geometric majorant gives
\[
 \left|T_\ell-5^{\ell/2}\sum_{n=1}^N
 \frac{x^{\ell n}+(-1)^\ell x^{2\ell n}}{(1-x^{3n})^\ell}\right|
 \leq\frac{2\cdot5^{\ell/2}r^{\ell(N+1)}}{(1-r^{3(N+1)})^\ell(1-r^\ell)} .
\]
At $N=600$ the right side is below $3.1\cdot10^{-83}$ for every index displayed. Values in the table are rounded to 15 decimal places. The sign of $T_1$ is determined by its defining series; the equation $T_1^2=T_2$ alone gives two possible square roots.
\begin{center}
\begin{tabular}{@{}lr@{}}\toprule
Expression & Numerical value\\\midrule
$T_1$ & $-2.479298354662861$\\
$T_2=T_1^2$ & $6.146920331433967$\\
$T_3$ & $-0.374629260249403$\\
$T_4$ & $6.966035075455954$\\
$T_5$ & $1.831243349153782$\\
$T_6$ & $10.355026974188011$\\
$T_7$ & $5.752405585353779$\\
$T_{12}$ & $54.082829641258435$\\
$T_{18}$ & $340.022309431742523$\\
$T_6T_{18}-T_{12}^2$ & $595.987723986014556$\\\bottomrule
\end{tabular}
\end{center}

\phantomsection\addcontentsline{toc}{section}{References}
\begin{thebibliography}{99}\small
\bibitem{AFdisjoint} B. Adamczewski and C. Faverjon, Linear disjointness and values of Mahler $M$-functions, preprint, 9 November 2026. \href{https://arxiv.org/abs/2611.12396}{arXiv:2611.12396}.

\bibitem{AFfinite} B. Adamczewski and C. Faverjon, Mahler's method in several variables and finite automata, \textit{Annals of Mathematics} \textbf{204} (2026), 455--533. \href{https://doi.org/10.4007/annals.2026.204.2.1}{doi:10.4007/annals.2026.204.2.1}. See also the addendum, pp. 535--543.

\bibitem{BBG} J.~M. Borwein, P.~B. Borwein, and F.~G. Garvan, Some cubic modular identities of Ramanujan, \emph{Trans. Amer. Math. Soc.} \textbf{343} (1994), 35--47. \href{https://doi.org/10.1090/S0002-9947-1994-1243610-6}{doi:10.1090/S0002-9947-1994-1243610-6}.

\bibitem{BG} L.~A. Borisov and P.~E. Gunnells, Toric modular forms of higher weight, \emph{J. Reine Angew. Math.} \textbf{560} (2003), 43--64. \href{https://doi.org/10.1515/crll.2003.060}{doi:10.1515/crll.2003.060}.

\bibitem{sn-sternA}
P.~Bundschuh, Transcendence and algebraic independence of series related
to Stern's sequence, \emph{Int. J. Number Theory} \textbf{8} (2012),
361--376. \href{https://doi.org/10.1142/S1793042112500212}{doi:10.1142/S1793042112500212}.

\bibitem{sn-stern}
P.~Bundschuh and K.~V\"a\"an\"anen, Algebraic independence of the generating
functions of Stern's sequence and of its twist, \emph{J. Th\'eor. Nombres
Bordeaux} \textbf{25} (2013), 43--57.
\href{https://doi.org/10.5802/jtnb.824}{doi:10.5802/jtnb.824}.

\bibitem{dsbv14} P. Bundschuh and K. V\"a\"an\"anen, Transcendence results on the generating functions of the characteristic functions of certain self-generating sets, \emph{Acta Arith.} \textbf{162} (2014), 273--288. \href{https://doi.org/10.4064/aa162-3-5}{doi:10.4064/aa162-3-5}.

\bibitem{sn-bvmahler}
P.~Bundschuh and K.~V\"a\"an\"anen, Algebraic independence of certain Mahler
functions and of their values, \emph{J. Aust. Math. Soc.} \textbf{98}
(2015), 289--310.
\href{https://doi.org/10.1017/S1446788714000524}{doi:10.1017/S1446788714000524}.

\bibitem{sn-lambert}
P.~Bundschuh and K.~V\"a\"an\"anen, Algebraic independence results on the
generating Lambert series of the powers of a fixed integer, \emph{Hardy--Ramanujan J.}
\textbf{38} (2015), 36--44.
\href{https://doi.org/10.46298/hrj.2015.1358}{doi:10.46298/hrj.2015.1358}.

\bibitem{dsbv15} P. Bundschuh and K. V\"a\"an\"anen, Transcendence results on the generating functions of the characteristic functions of certain self-generating sets, II, \emph{Acta Arith.} \textbf{167} (2015), 239--249. \href{https://doi.org/10.4064/aa167-3-2}{doi:10.4064/aa167-3-2}.

\bibitem{Chang} S. Chang, The Fibonacci sequence and an elliptic curve, \emph{J. Th\'eor. Nombres Bordeaux} \textbf{34} (2022), 483--495. \href{https://doi.org/10.5802/jtnb.1210}{doi:10.5802/jtnb.1210}.

\bibitem{Clery} F. Cl\'ery and G. van der Geer, Generating Picard modular forms by means of invariant theory, \S4. \href{https://arxiv.org/abs/2110.00849}{arXiv:2110.00849}.

\bibitem{ds09} K. Dilcher and K. B. Stolarsky, Stern polynomials and double-limit continued fractions, \emph{Acta Arith.} \textbf{140} (2009), 119--134. \href{https://doi.org/10.4064/aa140-2-2}{doi:10.4064/aa140-2-2}.

\bibitem{dsdhr} T. Dreyfus, C. Hardouin, and J. Roques, Hypertranscendence of solutions of Mahler equations, \emph{J. Eur. Math. Soc.} \textbf{20} (2018), 2209--2238. \href{https://doi.org/10.4171/JEMS/810}{doi:10.4171/JEMS/810}.

\bibitem{ESS} C. Elsner, S. Shimomura, and I. Shiokawa, Algebraic relations for reciprocal sums of Fibonacci numbers, \emph{Acta Arith.} \textbf{130} (2007), 37--60. \href{https://doi.org/10.4064/aa130-1-3}{doi:10.4064/aa130-1-3}.

\bibitem{EKS} I.~S. Eum, J.~K. Koo, and D.~H. Shin, On some applications of Eisenstein series, \emph{Publ. Math. Debrecen} \textbf{85} (2014). \href{https://doi.org/10.5486/PMD.2014.5813}{doi:10.5486/PMD.2014.5813}.

\bibitem{IdeGeometric} H. Ide, Algebraic independence results for a certain family of power series, infinite products, and Lambert type series, \textit{Osaka Journal of Mathematics} \textbf{60} (2023), 815--833. \href{https://doi.org/10.18910/93062}{doi:10.18910/93062}.

\bibitem{EulerStatus} R. Kalisch, T. Monderer, E. Leibtag, S. Weinbaum, M. Shalyt, and I. Kaminer, An infinite family of Ap\'ery-like recurrences for Euler's constant, preprint (2026). \href{https://arxiv.org/abs/2610.13919}{arXiv:2610.13919}.

\bibitem{Matsuda} K. Matsuda, Differential equations involving cubic theta functions and Eisenstein series, \emph{Osaka J. Math.} \textbf{57} (2020), 521--542. \href{https://doi.org/10.18910/76671}{doi:10.18910/76671}.

\bibitem{Nest} Yu.~V. Nesterenko, Modular functions and transcendence questions, \emph{Sb. Math.} \textbf{187} (1996), 1319--1348. \href{https://doi.org/10.1070/SM1996v187n09ABEH000158}{doi:10.1070/SM1996v187n09ABEH000158}.

\bibitem{sn-nish}
K.~Nishioka, \emph{Mahler Functions and Transcendence}, Lecture Notes in
Mathematics, vol.~1631, Springer, Berlin, 1996.
\href{https://doi.org/10.1007/BFb0093672}{doi:10.1007/BFb0093672}.

\bibitem{DLMF} F.~W.~J. Olver et al., eds., \emph{NIST Digital Library of Mathematical Functions}, \S5.8, infinite products for the gamma function. \href{https://dlmf.nist.gov/5.8}{NIST DLMF}.

\bibitem{TanakaGeom} T.-a. Tanaka, Algebraic independence of functions satisfying certain Mahler type functional equations and its applications, \textit{Seminar on Mathematical Sciences} \textbf{27}, Keio University (1999), 21--49. \href{https://www.math.keio.ac.jp/~takaaki/papers/FranJTanaka(manuscript).pdf}{Author's manuscript}.
\end{thebibliography}
\end{document}
